\documentclass[12pt]{article}
\usepackage{notes-common}
\usetikzlibrary{calc}
\geometry{margin=0.35in}
\AtBeginDocument{\fontsize{14}{17}\selectfont}

% Fri Oct 23 -- Magnetic fields + Biot-Savart (Knight 29.1-29.3). Same lesson
% for both tracks (vector algebra only, no calculus needed) -- matches
% circuits-intro-notes.tex's layout conventions (small margins, bigger base
% font). See biot-savart-intro-notes-key.tex for the filled-in version --
% same structure, wording and diagrams, every blank carries the answer.
% Built from Seth's rough notes (uploaded).
\newcommand{\eqline}[2][-0.5mm]{%
  \par\vspace{1mm}\noindent #2\par\vspace{#1}%
  \noindent\textcolor{rulehue}{\rule{\linewidth}{0.5pt}}\par\vspace{4mm}%
}
% A compact "(___, ___, ___)" vector-component blank, reused for every
% separation vector, unit vector and field vector below.
\newcommand{\vecbl}[1][1.3cm]{$\big($\,\bl[#1]\,, \bl[#1]\,, \bl[#1]\,$\big)$}

% ---- simple axis + point-charge diagram pieces, shared by all three figures ----
\newcommand{\axesxy}[1][2.6]{% small x-y axes through the origin, half-extent #1 (cm)
  \draw[-latex, line width=0.6pt, gray] (-#1,0) -- (#1,0) node[below right,scale=0.7]{$x$};
  \draw[-latex, line width=0.6pt, gray] (0,-#1) -- (0,#1) node[above left,scale=0.7]{$y$};
}
\newcommand{\qdot}[2]{% a source charge marker at (#1), labelled #2
  \filldraw[fill=accent!25, draw=accent] (#1) circle (2.6pt);
  \node[scale=0.75, anchor=north, yshift=-2pt] at (#1) {#2};
}
\newcommand{\pdot}[2]{% a field-point marker at (#1), labelled #2
  \filldraw[fill=white, draw=black] (#1) circle (2pt);
  \node[scale=0.75, anchor=south, yshift=2pt] at (#1) {#2};
}

\begin{document}

\noteshead{Intro to Biot--Savart}{}

\goal{
\begin{itemize}[itemsep=5pt,parsep=2pt,topsep=2pt]
\item Warm up with something you already know: superpose the $\vec E$ field
of a point charge at three locations, being explicit about separation
vectors and unit vectors.
\item Reuse that exact same geometry to introduce the Biot--Savart law for a
single \emph{moving} point charge, $\vec B$.
\item Compare $\vec E$ and $\vec B$ at those same three points --- same
source, same locations, very different field.
\item Scale up to three moving charges in a line and see their $\vec B$
contributions add (not cancel) at an off-axis point.
\item Make the leap from one moving charge ($q\vec v$) to a continuous
current ($I\,d\vec\ell$) --- the form of Biot--Savart we'll actually use for
wires, loops and solenoids.
\end{itemize}}

%============================================================
\nsec{Warm-up: $\vec E$ from a point charge}

A charge $q=+2.0\ \mu$C sits at $(-5,0,0)$ cm. Find $\vec E$ at three field
points: $P_1=(-9,3,0)$, $P_2=(-5,3,0)$, $P_3=(-1,3,0)$ (all in cm).

\begin{center}
\begin{tikzpicture}[x=0.5cm,y=0.5cm]
  \axesxy[5.5]
  \qdot{-5,0}{$q$}
  \pdot{-9,3}{$P_1$}
  \pdot{-5,3}{$P_2$}
  \pdot{-1,3}{$P_3$}
\end{tikzpicture}
\end{center}
\figtask{Sketch the three separation vectors $\vec r_1,\vec r_2,\vec r_3$ from
$q$ to each field point on the diagram above.}

\textbf{Step 1 --- separation vectors} $\vec r_i = \vec r_{P_i}-\vec r_q$ (cm):

$\vec r_1 = $ \vecbl

$\vec r_2 = $ \vecbl

$\vec r_3 = $ \vecbl

\textbf{Step 2 --- magnitudes and unit vectors} $\hat r_i = \vec r_i / r_i$:

$r_1 = $ \bl[1.2cm] cm, $\ \hat r_1 = $ \vecbl

$r_2 = $ \bl[1.2cm] cm, $\ \hat r_2 = $ \vecbl

$r_3 = $ \bl[1.2cm] cm, $\ \hat r_3 = $ \vecbl

\textbf{Step 3 --- Coulomb's law}, $\displaystyle \vec E = \frac{1}{4\pi
\varepsilon_0}\,\frac{q}{r^2}\,\hat r$:

$\vec E_1 = $ \vecbl[1.6cm] N/C

$\vec E_2 = $ \vecbl[1.6cm] N/C

$\vec E_3 = $ \vecbl[1.6cm] N/C

\checkpoint{$P_1$ and $P_3$ are the same distance from $q$. What's the same
about $\vec E_1$ and $\vec E_3$, and what's different? Why?}

%============================================================
\nsec{Biot--Savart: the same charge, now moving}

Same charge, same three points --- but now $q$ moves with velocity
$\vec v = (12,0,0)$ m/s. The Biot--Savart law for a single point charge is

\[
\vec B = \frac{\mu_0}{4\pi}\,\frac{q\,\vec v\times\hat r}{r^2}, \qquad
\frac{\mu_0}{4\pi} = 1\times10^{-7}\ \text{T}\!\cdot\!\text{m/A}.
\]

\textbf{Reuse your $\hat r_1,\hat r_2,\hat r_3$ from \S1} --- the geometry
hasn't changed, only the source. Compute $\vec v\times\hat r_i$ first, then
$\vec B_i$:

$\vec v\times\hat r_1 = $ \vecbl[1.4cm] m/s

$\vec v\times\hat r_2 = $ \vecbl[1.4cm] m/s

$\vec v\times\hat r_3 = $ \vecbl[1.4cm] m/s

$\vec B_1 = $ \vecbl[1.6cm] T

$\vec B_2 = $ \vecbl[1.6cm] T

$\vec B_3 = $ \vecbl[1.6cm] T

\checkpoint{Is $\vec B_1$ the mirror image of $\vec B_3$, the way $\vec E_1$
and $\vec E_3$ were? What's actually true about $\vec B_1$ and $\vec B_3$?}

%============================================================
\nsec{Compare $\vec E$ and $\vec B$}

\begin{itemize}[itemsep=6pt,parsep=2pt,topsep=2pt]
\item Which field is bigger here, $\vec E$ or $\vec B$ --- by roughly how
many orders of magnitude? \bl[9cm]
\item $\vec E_1,\vec E_2,\vec E_3$ point in \bl[7cm] directions (how do they
relate to $\hat r$?).
\item $\vec B_1,\vec B_2,\vec B_3$ all point \bl[7cm] (what direction, and
why does every one of them end up that way)?
\end{itemize}

\checkpoint{$\vec v$ points along $+\hat x$ and every field point here has
$z=0$. Could $\vec B$ ever have come out with a nonzero $x$- or
$y$-component at these points? Why or why not?}

%============================================================
\nsec{Three moving charges in a line}

Three charges, each $q=+1.0\ \mu$C, sit at $(1,-4,0)$, $(1,0,0)$, $(1,4,0)$
cm, each moving with $\vec v = (0,10,0)$ m/s. Find the total $\vec B$ at
field point $P=(-2,0,0)$ cm --- off to the side, not on the line of charges.

\begin{center}
\begin{tikzpicture}[x=0.45cm,y=0.45cm]
  \axesxy[6]
  \qdot{1,-4}{$q_1$}
  \qdot{1,0}{$q_2$}
  \qdot{1,4}{$q_3$}
  \pdot{-2,0}{$P$}
\end{tikzpicture}
\end{center}
\figtask{Sketch the three separation vectors from each charge to $P$.}

$\vec r_1 = $ \vecbl \hspace{0.3cm} $r_1 = $ \bl[1.1cm] cm \hspace{0.3cm}
$\hat r_1 = $ \vecbl

$\vec r_2 = $ \vecbl \hspace{0.3cm} $r_2 = $ \bl[1.1cm] cm \hspace{0.3cm}
$\hat r_2 = $ \vecbl

$\vec r_3 = $ \vecbl \hspace{0.3cm} $r_3 = $ \bl[1.1cm] cm \hspace{0.3cm}
$\hat r_3 = $ \vecbl

\textbf{Each charge's contribution} $\vec B_i = \dfrac{\mu_0}{4\pi}\dfrac{q\,
\vec v\times\hat r_i}{r_i^2}$:

$\vec B_1 = $ \vecbl[1.5cm] T \qquad $\vec B_2 = $ \vecbl[1.5cm] T

$\vec B_3 = $ \vecbl[1.5cm] T

\result[1.3cm]{$\vec B_{\text{total}} = \vec B_1+\vec B_2+\vec B_3 = $}

\checkpoint{Did the three contributions partly cancel, or did they add? Compare
to what you'd expect from three point \emph{charges'} $\vec E$ fields at an
off-axis point --- is $\vec B$ behaving the same way?}

%============================================================
\nsec{From one moving charge to a current}

A current is just a huge number of moving charges, confined to a wire. Take
a short piece of wire of length $d\ell$, carrying current $I$: in time $dt$,
the charge that passes through it is $dq$, moving at the drift velocity
$\vec v$ along the wire. The replacement rule is

\[
q\,\vec v \ \longrightarrow\ \bl[3cm]
\]

so the Biot--Savart law for one point charge, $\vec B = \dfrac{\mu_0}{4\pi}
\dfrac{q\,\vec v\times\hat r}{r^2}$, becomes the law for one current
element:

\result[2.4cm]{$d\vec B = $}

\checkpoint{Everything you did in \S1--\S4 above with separation vectors and
unit vectors --- does any of it change once $q\vec v$ becomes $I\,d\vec
\ell$? What's actually new?}

\end{document}
